\scanpage{169}\mspage{114}
foncteur $\inserted{X\longmapsto}\underline{\mathrm{Ch}}(X)$ (``chemins dans $X$'') sur $\mathrm{Cat}$, avec les structures de source, but et composition pour les chemins, et des chemins triviaux (agissant comme unités bilatères pour la composition). La construction d'un tel foncteur $C$ n'a pas l'air trivial et semble ne pas être connue -- j'ai peiné pour y arriver ! (\unsure{Avec|Avoir}, en plus, l'associativité de la composition \ldots) Ou bien utiliser, à la Quillen-\deleted{Thomason}, l'existence d'une classe de morphismes dits Q-fibrants, avec des bonnes propriétés de factorisation.
Ce n'est pas non plus évident \unsure{du|de|dans} $\mathrm{Cat}$, et c'est établi par Thomason, par une voie un peu compliquée. Peut-être serai-je amené à revenir dessus -- mais seulement en cas de besoin !

Par contre les \unsure{assertions} concernant les catégories $\underline{\mathrm{Fib}}\,S$, $\underline{\mathrm{Fibsc}}\,S$, ou $\underline{\mathrm{Cofib}}\,S$, $\underline{\mathrm{Cofibsc}}\,S$, et leurs variantes à indices $0$ (donc les huit catégories-\unsure{ailes|astres} du diagramme), et leurs relations à $\mathrm{Cat}_{/S}$ d'une part, à $\underline{\unsure{\mathrm{Fibhot}|\mathrm{Fiblat}}}\,S$ de l'autre (pour les cat. à indice $0$), ne peuvent être contenues dans un topo général sur des \deleted{\unsure{localiseurs}} cat. de modèles $(\unsure{M|\mathcal M},W)$, \inserted{me semble-t-il,} vu la nature très spéciale des quatre catégories désignées ci-dessus, de leur rôle très particulier dans le \deleted{\unsure{théorème}} formalisme catégorique. Or tout au moins les catégories ``scindées'' $\underline{\mathrm{Fibsc}}\,S$, $\underline{\mathrm{Cofibsc}}\,S$, et par suite aussi leurs variantes à indice $0$, jouent un rôle essentiel pour la théorie de $\mathrm{HOT}(S)$, vu que c'est elles qui
% archive sheet number: 168
